\documentclass[12pt,a4paper]{article}
\usepackage[utf8]{inputenc}
\usepackage[T2A]{fontenc}
\usepackage[russian]{babel}
\usepackage{amsmath}
\usepackage{amsfonts}
\usepackage{amssymb}
\usepackage{tikz}
\usepackage{circuitikz}
\usepackage{geometry}
\usepackage{graphicx}

\geometry{left=2.5cm,right=2.5cm,top=2.5cm,bottom=2.5cm}

\title{Тепловые флуктуации в интегральных схемах:\\
обобщение метода Найквиста для множественных элементов}
\author{}
\date{}

\begin{document}

\maketitle

\section{Введение: подход Найквиста}

Найквист (1928) рассмотрел два резистора $R_1$ и $R_2$, соединенных последовательно, каждый из которых генерирует тепловую ЭДС $E_1$ и $E_2$ соответственно.

\begin{figure}[h]
\centering
\begin{circuitikz}[scale=0.8]
\draw (0,0) to[R=$R_1$, v=$E_1$] (0,3) 
      to[short] (3,3) 
      to[R=$R_2$, v=$E_2$] (3,0) -- (0,0);
\end{circuitikz}
\caption{Два резистора с тепловыми ЭДС}
\end{figure}

Ток в цепи:
\begin{equation}
I = \frac{E_1}{R_1 + R_2}
\end{equation}

Мощность, рассеиваемая на $R_2$:
\begin{equation}
Q_2 = I^2 R_2 = \left(\frac{E_1}{R_1+R_2}\right)^2 R_2
\end{equation}

При $R_1 = R_2 = R$:
\begin{equation}
W = Q_2 = \left(\frac{E_1}{2R}\right)^2 R = \frac{E_1^2}{4R}
\end{equation}

\section{Спектральный подход Найквиста}

Рассмотрим линию передачи длиной $l$, соединяющую два сопротивления $R$.

\begin{figure}[h]
\centering
\begin{circuitikz}[scale=0.8]
\draw (0,0) to[R=$R$] (0,2) 
      to[transmission line, l=$l$] (4,2)
      to[R=$R$] (4,0) -- (0,0);
\end{circuitikz}
\caption{Линия передачи между сопротивлениями}
\end{figure}

Частоты собственных мод:
\begin{equation}
\nu_i = i \frac{v}{2l}, \quad i = 1, 2, 3, \ldots
\end{equation}

где $v$ --- скорость распространения волны.

Число мод в интервале частот $d\nu$:
\begin{equation}
N_{dof} = \frac{2l}{v} d\nu
\end{equation}

\subsection{Энергия мод}

Согласно классической статистической механике, на каждую степень свободы приходится энергия $kT$:
\begin{equation}
E_{\text{classical}} = \sum_i kT
\end{equation}

Согласно квантовой теории (Планк):
\begin{equation}
E_{\text{quantum}} = \sum_i \frac{h\nu_i}{e^{h\nu_i/kT} - 1}
\end{equation}

Полная энергия в интервале частот $d\nu$:
\begin{equation}
dW = N_{dof} \cdot \langle \varepsilon \rangle = \frac{2l}{v} d\nu \cdot \frac{h\nu}{e^{h\nu/kT} - 1}
\end{equation}

\section{Обобщение на RC-цепи}

Рассмотрим цепь с сопротивлением $R$ и емкостью $C$.

\begin{figure}[h]
\centering
\begin{circuitikz}[scale=0.8]
\draw (0,0) to[R=$R$] (0,2) 
      to[C=$C$] (2,2)
      to[short] (2,0) -- (0,0);
\end{circuitikz}
\caption{RC-цепь}
\end{figure}

Импеданс емкости:
\begin{equation}
Z_C = \frac{1}{j\omega C}
\end{equation}

Полный импеданс:
\begin{equation}
Z = R + \frac{1}{j\omega C}
\end{equation}

Ток:
\begin{equation}
I = \frac{E}{R + \frac{1}{j\omega C}}
\end{equation}

Мощность на резисторе:
\begin{equation}
I^2 R = C U^2 \frac{h\nu}{e^{h\nu/kT} + 1} d\nu
\end{equation}

\section{Нелинейные элементы: варикапы}

Для варикапа с нелинейной вольт-фарадной характеристикой $C(V)$:

\begin{equation}
Z_C(V) = \frac{1}{j\omega C(V)}
\end{equation}

где $C(V)$ --- функция напряжения.

\section{Системы из множества элементов}

\subsection{Цепочка из N одинаковых элементов}

\begin{figure}[h]
\centering
\begin{circuitikz}[scale=0.6]
\foreach \i in {0,1,2,3,4} {
    \draw (\i*2,0) to[R=$R$, l=$\ell$] (\i*2,1.5);
}
\draw (0,0) -- (8,0);
\draw (0,1.5) -- (8,1.5);
\end{circuitikz}
\caption{Цепочка из 5 элементов}
\end{figure}

Для цепочки из $N$ элементов длиной $\ell$ каждый:

Частоты мод:
\begin{equation}
\nu_i = i \frac{v}{2N\ell}, \quad i = 1, 2, \ldots
\end{equation}

Длины волн:
\begin{equation}
\lambda_i = \frac{2N\ell}{i}
\end{equation}

\subsection{Случай N = 5}

\begin{equation}
\nu_i = i \frac{v}{10\ell}, \quad i = 1, 2, 3, 4, 6, 7, 8, 9, \ldots
\end{equation}
(исключая $i$, делящиеся на 5)

\begin{equation}
\lambda_i = \frac{10\ell}{i}
\end{equation}

Плотность состояний:
\begin{equation}
\frac{di_{i5}}{d\nu} = \frac{10\ell}{v} \cdot \frac{4}{5} = \frac{8\ell}{v}
\end{equation}

\subsection{Случай N = 7}

\begin{equation}
\nu_i = i \frac{v}{12\ell}, \quad i = 1, 5, 7, 11, 13, 17, 19, 23, \ldots
\end{equation}
(исключая $i$, делящиеся на 2, 3, 4, 6)

\begin{equation}
\lambda_i = \frac{12\ell}{i}
\end{equation}

\section{Плотность состояний для различных конфигураций}

\subsection{Два элемента}

\begin{equation}
\nu_i = i \frac{v}{2\ell}, \quad i = 1, 2, 3, 4
\end{equation}

\begin{equation}
\lambda_i = \frac{2\ell}{i}
\end{equation}

\begin{equation}
\frac{di}{d\nu} = \frac{2\ell}{v}
\end{equation}

\subsection{Три элемента}

\begin{equation}
\nu_j = j \frac{v}{4\ell}, \quad j = 1, 3, 5, 7, \ldots
\end{equation}
(нечетные $j$)

\begin{equation}
\lambda_j = \frac{4\ell}{j}
\end{equation}

Плотность состояний:
\begin{equation}
N_{dof} = \left(\frac{di}{d\nu} + \frac{dj}{d\nu}\right) d\nu = \left(\frac{2\ell}{v} + \frac{2\ell}{v}\right) d\nu = \frac{4\ell}{v} d\nu
\end{equation}

\subsection{Четыре элемента}

\begin{equation}
\nu_k = k \frac{v}{6\ell}, \quad k = 1, 2, 4, 5, 7, 8, 10, 11, 13, 14, \ldots
\end{equation}
(исключая $k$, делящиеся на 3)

Плотность состояний:
\begin{equation}
N_{dof} = \left(\frac{di_2}{d\nu} + 2\frac{di_3}{d\nu} + \frac{dk_4}{d\nu}\right) d\nu = \left(\frac{2\ell}{v} + 2\cdot\frac{2\ell}{v} + \frac{4\ell}{v}\right) d\nu = \frac{10\ell}{v} d\nu
\end{equation}

\section{Обобщенная формула для N элементов}

Для цепочки из $N$ элементов:

\begin{equation}
\nu_i = i \frac{v}{2N\ell}
\end{equation}

\begin{equation}
\lambda_i = \frac{2N\ell}{i}
\end{equation}

Плотность состояний:
\begin{equation}
N_{dof} = \sum_{m} c_m \frac{d i_m}{d\nu}
\end{equation}

где $c_m$ --- коэффициенты, зависящие от конфигурации.

\section{Применение к интегральным схемам}

Рассмотрим интегральную схему, состоящую из $N$ параллельно соединенных элементов.

\subsection{Суммирование мощностей}

Если каждый элемент генерирует мощность теплового шума $P_i$, то полная мощность:

\begin{equation}
P_{\text{total}} = \sum_{i=1}^{N} P_i
\end{equation}

При одинаковых элементах:
\begin{equation}
P_{\text{total}} = N \cdot P_{\text{single}}
\end{equation}

\subsection{Плотность размещения}

На площади $S = 1 \text{ см}^2$ можно разместить $N$ элементов размером $a \times b$:

\begin{equation}
N = \frac{S}{a \cdot b} = \frac{10^{-4} \text{ м}^2}{(35 \times 10^{-9})(30 \times 10^{-9}) \text{ м}^2} \approx 9.5 \times 10^{10}
\end{equation}

\section{Выводы}

1. Метод Найквиста может быть обобщен на системы из множества элементов.

2. Плотность состояний зависит от конфигурации соединений.

3. Для интегральных схем возможно суммирование мощностей тепловых флуктуаций.

4. Нелинейные элементы (варикапы) требуют учета зависимости $C(V)$.

\end{document}